\documentclass[11pt,a4paper]{article}

\usepackage[T1]{fontenc}
\usepackage[utf8]{inputenc}
\usepackage[ngerman]{babel}
\usepackage{lmodern}
\usepackage{amsmath,amssymb,amsthm,mathtools}
\usepackage{geometry}
\usepackage{booktabs}
\usepackage{hyperref}
\usepackage{microtype}
\usepackage{enumitem}

\geometry{margin=2.7cm}

\title{Lagrange-Vierlinge im Primzahlraum \(p>3\)\\
	und zyklische Schließbarkeit im Sehnenmodell}
\author{Thomas Hoffbauer}
\date{\today}

\newtheorem{satz}{Satz}
\newtheorem{definition}{Definition}
\newtheorem{lemma}{Lemma}
\newtheorem{korollar}{Korollar}
\theoremstyle{remark}
\newtheorem{bemerkung}{Bemerkung}
\newtheorem{beispiel}{Beispiel}

\begin{document}
	\maketitle
	
	\begin{abstract}
		Wir untersuchen Vierquadrate-Zerlegungen von Primzahlen \(p>3\) und interpretieren sie als \emph{Lagrange-Vierlinge} im Primzahlraum. Jede solche Darstellung
		\[
		p = E^2 + A^2 + B^2 + C^2
		\]
		wird als Viererstruktur aufgefasst, deren Komponenten in ein zyklisches Sehnenmodell eingebettet werden. Dies führt zu einer zusätzlichen geometrischen Schließungsbedingung, die innerhalb der arithmetisch durch Lagrange garantierten Darstellungen selektiv wirkt.
		
		Für die ersten 1000 Primzahlen wurde numerisch untersucht, welche Lagrange-Vierlinge im harten Modell \(l_i = E_i^2\) beziehungsweise im weichen Modell \(l_i = E_i\) zyklisch schließbar sind. Dabei zeigt sich, dass einrastfähige Vierlinge im Mittel deutlich höhere Entropie, geringere Konzentration und größere Balance aufweisen als nicht-einrastfähige. Zugleich verschärft die quadratische Gewichtung die Schließungsbedingung spürbar: Im harten Modell besitzen 983 der ersten 1000 Primzahlen mindestens einen fit-Vierling, im weichen Modell 993.
		
		Damit erscheint der Primzahlraum \(p>3\) nicht nur als Menge arithmetisch darstellbarer Objekte, sondern als Raum innerer Vierer-Geometrien mit unterschiedlich ausgeprägter zyklischer Verträglichkeit.
	\end{abstract}
	
	\section{Einleitung}
	
	Der Vier-Quadrate-Satz von Lagrange garantiert, dass jede natürliche Zahl als Summe von vier Quadraten darstellbar ist. Für Primzahlen \(p>3\) ergibt sich damit eine universelle Viererstruktur
	\[
	p = E^2 + A^2 + B^2 + C^2.
	\]
	Im Folgenden nennen wir jede solche Darstellung einen \emph{Lagrange-Vierling} von \(p\).
	
	Die klassische Aussage Lagranges ist rein existenziell: Sie sichert das Vorhandensein einer Vierquadrate-Zerlegung, sagt aber zunächst nichts über die innere Gestalt dieser Zerlegung aus. Die leitende Idee der vorliegenden Arbeit besteht darin, Lagrange-Vierlinge nicht nur als additive Repräsentationen, sondern als geometrisch interpretierbare Vierer-Konfigurationen aufzufassen.
	
	Dazu werden die vier Beiträge in ein zyklisches Sehnenmodell eingebettet. Hieraus entsteht eine zusätzliche Schließungsbedingung: Nur ein Teil der arithmetisch möglichen Lagrange-Vierlinge ist geometrisch kompatibel. Auf diese Weise wird der Primzahlraum \(p>3\) zu einem Raum innerer Vierer-Geometrien, in dem zwischen schließbaren, beinahe schließbaren und deutlich nicht-schließbaren Konfigurationen unterschieden werden kann.
	
	Der konzeptionelle Ausgangspunkt lautet daher:
	\begin{quote}
		Lagrange liefert die arithmetische Universalität, das Sehnenmodell führt eine geometrische Selektivität ein.
	\end{quote}
	
	Numerisch zeigt sich, dass diese Selektivität keineswegs zufällig ist. Vielmehr korreliert zyklische Schließbarkeit stark mit einer balancierten Verteilung der inneren Lasten eines Lagrange-Vierlings. Dadurch entsteht im Primzahlraum eine zusätzliche arithmetisch-geometrische Feinstruktur.
	
	\section{Lagrange-Vierlinge}
	
	\begin{definition}
		Sei \(p>3\) eine Primzahl. Ein \emph{Lagrange-Vierling} von \(p\) ist ein Tupel
		\[
		(E,A,B,C)\in \mathbb Z_{\ge 0}^4
		\]
		mit
		\[
		E \ge A \ge B \ge C
		\qquad\text{und}\qquad
		E^2 + A^2 + B^2 + C^2 = p.
		\]
		Die Menge aller Lagrange-Vierlinge von \(p\) bezeichnen wir mit
		\[
		\mathcal L(p).
		\]
	\end{definition}
	
	\begin{satz}[Lagrange]
		Für jede natürliche Zahl \(n\) existiert eine Darstellung
		\[
		n = x_1^2 + x_2^2 + x_3^2 + x_4^2
		\]
		mit \(x_i \in \mathbb Z\). Insbesondere gilt für jede Primzahl \(p>3\):
		\[
		\mathcal L(p)\neq\emptyset.
		\]
	\end{satz}
	
	\begin{bemerkung}
		Im numerischen Teil werden ungeordnete Repräsentanten
		\[
		E \ge A \ge B \ge C \ge 0
		\]
		verwendet. Für feinere zyklische Fragestellungen kann später auch die Menge aller verschiedenen Permutationen eines Lagrange-Vierlings betrachtet werden.
	\end{bemerkung}
	
	\section{Das zyklische Sehnenmodell}
	
	\subsection{Zwei Gewichtungsregime}
	
	Zu einem Lagrange-Vierling \((E,A,B,C)\in\mathcal L(p)\) ordnen wir vier Längen \(l_1,l_2,l_3,l_4\) zu. Wir betrachten zwei Varianten:
	
	\begin{enumerate}[label=\arabic*)]
		\item \textbf{Harte Gewichtung}
		\[
		l_1 = E^2,\qquad l_2 = A^2,\qquad l_3 = B^2,\qquad l_4 = C^2.
		\]
		
		\item \textbf{Weiche Gewichtung}
		\[
		l_1 = E,\qquad l_2 = A,\qquad l_3 = B,\qquad l_4 = C.
		\]
	\end{enumerate}
	
	Die harte Gewichtung verstärkt große Komponenten überproportional, die weiche Gewichtung behandelt sie linear.
	
	\subsection{Schließungsbedingung}
	
	Für eine Sehne der Länge \(l\) in einem Kreis vom Radius \(r\) gilt
	\[
	\theta(l;r)=2\arcsin\!\left(\frac{l}{2r}\right).
	\]
	Ein Lagrange-Vierling heißt \emph{einrastfähig} oder kurz \emph{fit}, wenn es einen Radius \(r>0\) gibt mit
	\[
	\theta(l_1;r)+\theta(l_2;r)+\theta(l_3;r)+\theta(l_4;r)=2\pi.
	\]
	Äquivalent:
	\[
	\sum_{i=1}^4 \arcsin\!\left(\frac{l_i}{2r}\right)=\pi.
	\]
	
	\begin{definition}
		Existiert ein Radius \(r>0\), der die Gleichung
		\[
		\sum_{i=1}^4 \arcsin\!\left(\frac{l_i}{2r}\right)=\pi
		\]
		erfüllt, so nennen wir ihn die \emph{Einrastzahl} des Lagrange-Vierlings und schreiben
		\[
		\kappa(E,A,B,C)=r.
		\]
	\end{definition}
	
	\begin{lemma}
		Die Funktion
		\[
		f(r)=\sum_{i=1}^4 \arcsin\!\left(\frac{l_i}{2r}\right)
		\]
		ist auf ihrem Definitionsbereich streng monoton fallend.
	\end{lemma}
	
	\begin{proof}
		Jeder Summand
		\[
		r\mapsto \arcsin\!\left(\frac{l_i}{2r}\right)
		\]
		ist für \(r>\frac{l_i}{2}\) streng fallend, weil \(r\mapsto \frac{l_i}{2r}\) streng fällt und \(\arcsin\) auf \((0,1)\) streng wächst. Die Summe streng fallender Funktionen ist wieder streng fallend.
	\end{proof}
	
	\begin{korollar}
		Falls \(\kappa(E,A,B,C)\) existiert, dann ist sie eindeutig.
	\end{korollar}
	
	\section{Kritischer Defekt und Defektgeometrie}
	
	\begin{definition}
		Sei
		\[
		r_{\min}=\frac12\max(l_1,l_2,l_3,l_4).
		\]
		Dann definieren wir den \emph{kritischen Defekt} eines Lagrange-Vierlings durch
		\[
		\delta(E,A,B,C)=\pi-f(r_{\min}),
		\qquad
		f(r)=\sum_{i=1}^4 \arcsin\!\left(\frac{l_i}{2r}\right).
		\]
	\end{definition}
	
	Da \(f(r)\) streng fällt, ist \(f(r_{\min})\) der maximal erreichbare Wert. Daher gilt:
	
	\[
	\delta>0 \;\Longrightarrow\; \text{non-fit},
	\qquad
	\delta=0 \;\Longrightarrow\; \text{Grenzfall},
	\qquad
	\delta<0 \;\Longrightarrow\; \text{fit}.
	\]
	
	Numerisch wurde fit äquivalent durch \(f(r_{\min})\ge\pi\) mit Toleranz geprüft; fit-Fälle erscheinen dann mit Defekt \(0\).
	
	\begin{bemerkung}
		Der kritische Defekt erlaubt es, auch non-fit-Fälle auf einer kontinuierlichen Skala zu ordnen. Kleine positive Werte entsprechen beinahe schließbaren Konfigurationen, große Werte deutlich rigiden Konfigurationen.
	\end{bemerkung}
	
	\section{Verteilungsmaße eines Lagrange-Vierlings}
	
	Zur Beschreibung der inneren Lastverteilung verwenden wir die normierten Gewichte
	\[
	p_i=\frac{l_i}{l_1+l_2+l_3+l_4}.
	\]
	
	Daraus definieren wir:
	
	\begin{definition}
		Die \emph{Entropie} eines Lagrange-Vierlings sei
		\[
		H=-\sum_{i=1}^4 p_i\log p_i.
		\]
	\end{definition}
	
	\begin{definition}
		Die \emph{Konzentration} sei
		\[
		K=\sum_{i=1}^4 p_i^2.
		\]
	\end{definition}
	
	\begin{definition}
		Der \emph{Balanceindex} sei
		\[
		\operatorname{bal}
		=
		\frac{\min(l_1,l_2,l_3,l_4)}{\max(l_1,l_2,l_3,l_4)}.
		\]
	\end{definition}
	
	Hohe Entropie und hoher Balanceindex entsprechen einer relativ gleichmäßigen Lastverteilung; hohe Konzentration entspricht einer dominanten Hauptkomponente.
	
	\section{Numerische Untersuchung der ersten 1000 Primzahlen}
	
	Es wurden die ersten 1000 Primzahlen
	\[
	2,3,5,\dots,7919
	\]
	numerisch untersucht. Im Mittelpunkt stehen die Primzahlen \(p>3\); die kleinen Fälle \(2\) und \(3\) wurden der Vollständigkeit halber mitgeführt.
	
	\subsection{Harte Gewichtung \(l_i=E_i^2\)}
	
	Im harten Modell besitzen 983 der ersten 1000 Primzahlen mindestens einen fit-Lagrange-Vierling; 17 Primzahlen besitzen keinen fit-Vierling. Diese 17 Primzahlen sind:
	\[
	5,\;7,\;11,\;17,\;29,\;47,\;53,\;61,\;71,\;83,\;131,\;149,\;197,\;233,\;239,\;257,\;269.
	\]
	
	\subsection{Statistik über alle Darstellungen}
	
	Die statistische Auswertung über sämtliche untersuchten Lagrange-Vierlinge der ersten 1000 Primzahlen ergibt im harten Modell:
	
	\begin{center}
		\begin{tabular}{lcccc}
			\toprule
			Gruppe & Anzahl & mittlere Entropie \(H\) & mittlere Konzentration \(K\) & mittlere Balance \\
			\midrule
			fit     & 11046 & 1.1898 & 0.3262 & 0.1977 \\
			non-fit & 77685 & 0.8324 & 0.5302 & 0.0454 \\
			\bottomrule
		\end{tabular}
	\end{center}
	
	Damit zeigen fit-Vierlinge im Mittel:
	
	\begin{itemize}
		\item deutlich \emph{höhere Entropie},
		\item deutlich \emph{geringere Konzentration},
		\item deutlich \emph{größere Balance}.
	\end{itemize}
	
	\subsection{Interpretation}
	
	Die zyklische Schließbarkeit ist damit klar \emph{nicht} bloß eine Umformulierung des Vier-Quadrate-Satzes. Vielmehr wirkt das Sehnenmodell als geometrischer Filter innerhalb der von Lagrange garantierten Darstellungen. Es bevorzugt solche Lagrange-Vierlinge, bei denen die Last nicht zu stark auf eine dominante Hauptkomponente konzentriert ist.
	
	Der zentrale numerische Befund lautet daher:
	\begin{quote}
		Einrastfähigkeit korreliert stark mit einer ausgeglicheneren Lastverteilung des Lagrange-Vierlings.
	\end{quote}
	
	\section{Fast einrastende non-fit-Fälle}
	
	Unter den 17 non-fit-Primzahlen im harten Modell gibt es sehr unterschiedliche Defektgrößen. Einige Fälle sind deutlich non-fit, andere liegen extrem nahe an der Schließung.
	
	Besonders kleine Defekte wurden etwa für folgende Primzahlen gefunden:
	\[
	197,\qquad 61,\qquad 239,\qquad 233,\qquad 149.
	\]
	
	Ein markanter Grenzfall ist \(p=197\) mit dem Lagrange-Vierling
	\[
	(9,8,6,4),
	\]
	für den numerisch
	\[
	\delta \approx 3.93\times 10^{-4}
	\]
	gefunden wurde. Dies zeigt, dass die Klasse der non-fit-Vierlinge selbst eine feine innere Struktur besitzt: Sie reicht von stark rigiden bis zu nahezu schließbaren Konfigurationen.
	
	\section{Vergleich zwischen harter und weicher Gewichtung}
	
	Im weichen Modell
	\[
	l_i=E_i
	\]
	besitzen 993 der ersten 1000 Primzahlen mindestens einen fit-Vierling. Der Vergleich ergibt somit:
	
	\begin{center}
		\begin{tabular}{lc}
			\toprule
			Größe & Wert \\
			\midrule
			fit im harten Modell & 983 \\
			fit im weichen Modell & 993 \\
			Verbesserungen hard \(\to\) soft & 10 \\
			Verschlechterungen hard \(\to\) soft & 0 \\
			in beiden Modellen non-fit & 7 \\
			\bottomrule
		\end{tabular}
	\end{center}
	
	Dies zeigt eindeutig:
	\begin{quote}
		Die quadratische Gewichtung verschärft die geometrische Schließungsbedingung.
	\end{quote}
	
	Große Komponenten werden im harten Modell überproportional verstärkt; dadurch können Lagrange-Vierlinge, die unter linearer Gewichtung noch kompatibel sind, ihre zyklische Schließbarkeit verlieren.
	
	\section{Beste Darstellung pro Primzahl}
	
	Für jede Primzahl wurde zusätzlich eine \emph{beste Darstellung} ausgewählt:
	
	\begin{itemize}
		\item falls fit-Vierlinge existieren: der fit-Vierling mit kleinster Einrastzahl \(\kappa\),
		\item sonst: der non-fit-Vierling mit kleinstem Defekt \(\delta\).
	\end{itemize}
	
	Diese Auswahl liefert für jede Primzahl einen kanonischen Repräsentanten ihrer inneren Geometrie im jeweiligen Modell. Im harten Modell zeigt sich dabei häufig, dass die besten fit-Vierlinge eine bemerkenswert gleichmäßige Struktur haben, etwa:
	\[
	(3,3,3,2),\quad
	(4,3,3,3),\quad
	(5,4,4,4),\quad
	(7,6,6,6).
	\]
	Dies stützt erneut die Vermutung, dass Schließbarkeit eng mit einer balancierten Lastverteilung zusammenhängt.
	
	Zugleich zeigt der Vergleich zwischen harter und weicher Gewichtung, dass sich nicht nur die Fit-Quote, sondern teilweise auch der optimale Vierling einer Primzahl ändert. Das Modell reagiert also empfindlich auf das gewählte Gewichtungsregime.
	
	\section{Schlussfolgerung}
	
	Die vorliegende Untersuchung zeigt, dass Vierquadrate-Zerlegungen von Primzahlen \(p>3\) durch ein zyklisches Sehnenmodell geometrisch verfeinert werden können. Die durch Lagrange garantierte Existenz von Vierquadrate-Zerlegungen bildet dabei die arithmetische Basis; die zusätzliche Kreis-Schließungsbedingung führt eine nichttriviale geometrische Selektivität ein.
	
	Die numerischen Resultate legen nahe:
	
	\begin{enumerate}[label=\arabic*)]
		\item Einrastfähigkeit ist stark mit hoher Entropie und Balance sowie geringer Konzentration korreliert.
		\item Die harte Gewichtung \(l_i = E_i^2\) wirkt als Verstärker struktureller Ungleichverteilung.
		\item Nicht-einrastende Primzahlen bilden keine homogene Klasse, sondern reichen von stark rigiden bis zu nahezu kritischen Grenzfällen.
		\item Die Wahl des Gewichtungsregimes verändert nicht nur die Fit-Quote, sondern teilweise auch die optimale Darstellung einer Primzahl.
	\end{enumerate}
	
	Damit entsteht aus der klassischen Vierquadrate-Theorie eine arithmetisch-geometrische Verfeinerung, in der nicht nur die Existenz von Darstellungen zählt, sondern auch deren innere zyklische Verträglichkeit.
	
	\section*{Ausblick}
	
	Naheliegende Fortsetzungen dieser Arbeit sind:
	
	\begin{itemize}
		\item Untersuchung geordneter statt nur ungeordneter Lagrange-Vierlinge,
		\item Vergleich der Klassen \(p\equiv 1 \pmod 4\) und \(p\equiv 3 \pmod 4\),
		\item asymptotische Analyse der Fit-Quote für größere Primzahlbereiche,
		\item Untersuchung anderer Gewichtungsregime \(l_i = E_i^\alpha\),
		\item quaternionische Interpretation der Viererstruktur über
		\[
		N(E+Ai+Bj+Ck)=E^2+A^2+B^2+C^2.
		\]
	\end{itemize}
	
	\section{Diskrete Pfade, Transferoperator und BM-Variationsprinzip}
	
	\subsection{Motivation}
	
	Die formale Feynman-Darstellung eines Propagators als Summe über alle Pfade
	\[
	K(x,y;T)=\int_{\gamma(0)=x}^{\gamma(T)=y}
	\exp\!\left(\frac{i}{\hbar}S[\gamma]\right)\,\mathcal D\gamma
	\]
	ist konzeptionell außerordentlich fruchtbar, aber analytisch oft schwer kontrollierbar. 
	Im Rahmen des Bamberger Modells liegt daher die Idee nahe, das Pfadintegral nicht als Summe über beliebige kontinuierliche Kurven, sondern als \emph{diskrete Zustandssumme über strukturierte Pfade} zu formulieren.
	
	Der leitende Gedanke lautet:
	\begin{quote}
		Nicht jeder formal mögliche Pfad trägt im gleichen Maße zur Dynamik bei. 
		Vielmehr sollen solche Pfade bevorzugt werden, die im Sinne des Modells strukturell kohärent und defektarm sind.
	\end{quote}
	
	Damit tritt an die Stelle der Summation über beliebige Bahnen eine Summation über diskrete Historien im Zustandsraum des Modells, gewichtet durch Phase und Defekt.
	
	\subsection{BM-Zustandsraum}
	
	\begin{definition}
		Ein \emph{BM-Zustand} ist ein Element eines diskreten Zustandsraums
		\[
		\mathcal S.
		\]
		In einer minimalen Version kann ein Zustand bereits durch einen Lagrange-Vierling
		\[
		X=(E,A,B,C)
		\]
		repräsentiert werden. 
		In reicheren Versionen können zusätzliche Komponenten wie Modulo-Klassen, Shell-Lagen, Phasenmarken oder quaternionische Orientierungen hinzutreten, etwa
		\[
		X=(E,A,B,C;\sigma,\mu,\nu,\dots).
		\]
	\end{definition}
	
	\begin{bemerkung}
		Für den hier skizzierten Ansatz genügt es, \(\mathcal S\) als diskrete Menge strukturierter Zustände aufzufassen. 
		Die konkrete Wahl der Zustandskoordinaten kann je nach Ausbaustufe des Bamberger Modells verfeinert werden.
	\end{bemerkung}
	
	\subsection{Erlaubte Übergänge und diskrete Pfade}
	
	\begin{definition}
		Es sei
		\[
		\rightsquigarrow \;\subseteq\; \mathcal S\times\mathcal S
		\]
		eine Übergangsrelation auf \(\mathcal S\). 
		Ein Ausdruck
		\[
		X\rightsquigarrow Y
		\]
		bedeutet, dass der Übergang von \(X\) nach \(Y\) im Sinne des Modells erlaubt ist.
	\end{definition}
	
	\begin{definition}
		Ein \emph{diskreter BM-Pfad} der Länge \(N\) von \(X\) nach \(Y\) ist eine Folge
		\[
		\gamma=(X_0,X_1,\dots,X_N)
		\]
		mit
		\[
		X_0=X,\qquad X_N=Y,
		\qquad X_k\rightsquigarrow X_{k+1}\quad (k=0,\dots,N-1).
		\]
		Die Menge aller solcher Pfade sei
		\[
		\Gamma_N(X,Y).
		\]
	\end{definition}
	
	Damit wird der Pfadraum des Modells nicht durch willkürliche Kurven, sondern durch Wege in einem diskreten Strukturgraphen beschrieben.
	
	\subsection{Lokaler Defekt und lokale Phase}
	
	Um Beiträge einzelner Pfade zu gewichten, werden zwei lokale Größen eingeführt: ein Defekt und eine Phase.
	
	\begin{definition}
		Zu jedem erlaubten Übergang
		\[
		X\rightsquigarrow Y
		\]
		gehöre ein \emph{lokaler Defekt}
		\[
		d(X,Y)\ge 0
		\]
		und eine \emph{lokale Phase}
		\[
		\phi(X,Y)\in\mathbb R.
		\]
	\end{definition}
	
	Der Defekt misst, wie stark ein Übergang gegen die innere Struktur des Modells verstößt; die Phase kodiert den oszillatorischen Anteil der Dynamik.
	
	In einer ersten Ausbaustufe kann der Defekt aus einem lokalen Vierlingsdefekt, aus Modulo-Inkompatibilitäten, aus Shell-Wechseln oder aus topologischen Nichtschließungen zusammengesetzt werden, etwa in der Form
	\[
	d(X,Y)
	=
	\alpha\, d_{\mathrm{quad}}(X,Y)
	+\beta\, d_{\mathrm{mod}}(X,Y)
	+\gamma\, d_{\mathrm{top}}(X,Y).
	\]
	
	\subsection{Lagrange-Vierlinge als lokales Defektmodul}
	
	Ist ein Zustand \(X\) durch einen Lagrange-Vierling
	\[
	X=(E,A,B,C)
	\]
	gegeben, so kann der zuvor eingeführte kritische Defekt
	\[
	\delta(E,A,B,C)=\pi-f(r_{\min})
	\]
	mit
	\[
	f(r)=\sum_{i=1}^4 \arcsin\!\left(\frac{l_i}{2r}\right)
	\]
	als lokales Strukturmaß benutzt werden.
	
	Eine einfache Wahl ist
	\[
	d_{\mathrm{quad}}(X)=\max(0,\delta(E,A,B,C)).
	\]
	Alternativ kann auch der Übergang selbst mit einem Defekt bewertet werden, etwa durch
	\[
	d_{\mathrm{quad}}(X,Y)
	=
	d_{\mathrm{quad}}(Y)
	+\eta\,|d_{\mathrm{quad}}(Y)-d_{\mathrm{quad}}(X)|.
	\]
	
	Damit gehen die zuvor untersuchten Lagrange-Vierlinge unmittelbar als lokales Mikromodul in die Pfadbeschreibung ein.
	
	\subsection{Diskrete BM-Wirkung}
	
	\begin{definition}
		Die \emph{diskrete BM-Lagrangefunktion} sei eine Funktion
		\[
		L_{\mathrm{BM}}:\mathcal S\times\mathcal S\to\mathbb C
		\]
		der Form
		\[
		L_{\mathrm{BM}}(X,Y)
		=
		\phi(X,Y)+ i\lambda\, d(X,Y),
		\qquad \lambda>0.
		\]
	\end{definition}
	
	\begin{definition}
		Für einen diskreten Pfad
		\[
		\gamma=(X_0,\dots,X_N)
		\]
		definieren wir die \emph{BM-Wirkung}
		\[
		\mathcal S_{\mathrm{BM}}[\gamma]
		=
		\sum_{k=0}^{N-1}L_{\mathrm{BM}}(X_k,X_{k+1}).
		\]
		Schreibt man
		\[
		\Phi[\gamma]
		=
		\sum_{k=0}^{N-1}\phi(X_k,X_{k+1}),
		\qquad
		\Delta[\gamma]
		=
		\sum_{k=0}^{N-1}d(X_k,X_{k+1}),
		\]
		so gilt
		\[
		\mathcal S_{\mathrm{BM}}[\gamma]
		=
		\Phi[\gamma]+ i\lambda\,\Delta[\gamma].
		\]
	\end{definition}
	
	Daraus ergibt sich als Pfadamplitude
	\[
	\mathcal A(\gamma)
	=
	\exp\!\left(i\mathcal S_{\mathrm{BM}}[\gamma]\right)
	=
	\exp\!\left(i\Phi[\gamma]-\lambda\Delta[\gamma]\right).
	\]
	
	\begin{bemerkung}
		Die Phase \(\Phi[\gamma]\) erzeugt Interferenz, der Defekt \(\Delta[\gamma]\) dämpft strukturell unverträgliche Pfade. 
		Das Modell verbindet damit eine oszillatorische mit einer selektiven Komponente.
	\end{bemerkung}
	
	\subsection{Diskreter BM-Propagator}
	
	\begin{definition}
		Für \(X,Y\in\mathcal S\) definieren wir den BM-Propagator der Pfadlänge \(N\) durch
		\[
		K_N(X,Y)
		=
		\sum_{\gamma\in\Gamma_N(X,Y)}
		\mathcal A(\gamma)
		=
		\sum_{\gamma\in\Gamma_N(X,Y)}
		\exp\!\left(i\Phi[\gamma]-\lambda\Delta[\gamma]\right).
		\]
	\end{definition}
	
	Formal ergibt sich daraus der volle Propagator
	\[
	K(X,Y)
	=
	\sum_{N\ge 0}K_N(X,Y),
	\]
	gegebenenfalls mit zusätzlicher Längenstrafe \(\alpha^N\), \(0<\alpha\le 1\).
	
	\subsection{Transferoperator}
	
	\begin{definition}
		Der \emph{BM-Transferoperator} \(T\) sei durch seine Matrixelemente
		\[
		T_{XY}
		=
		\chi(X\rightsquigarrow Y)\,
		\exp\!\left(i\phi(X,Y)-\lambda d(X,Y)\right)
		\]
		gegeben, wobei \(\chi(X\rightsquigarrow Y)\in\{0,1\}\) die Übergangsrelation kodiert.
	\end{definition}
	
	\begin{satz}
		Für jedes \(N\in\mathbb N\) gilt
		\[
		(T^N)_{XY}
		=
		\sum_{\gamma\in\Gamma_N(X,Y)}
		\exp\!\left(i\Phi[\gamma]-\lambda\Delta[\gamma]\right).
		\]
		Insbesondere ist
		\[
		(T^N)_{XY}=K_N(X,Y).
		\]
	\end{satz}
	
	\begin{proof}
		Für \(N=1\) ist die Aussage genau die Definition von \(T\). 
		Der Induktionsschritt folgt aus
		\[
		(T^{N+1})_{XY}
		=
		\sum_{Z\in\mathcal S}(T^N)_{XZ}T_{ZY},
		\]
		also aus der Zerlegung jedes \((N+1)\)-Schritt-Pfades in einen \(N\)-Schritt-Pfad von \(X\) nach \(Z\) und einen letzten Schritt von \(Z\) nach \(Y\).
	\end{proof}
	
	\begin{bemerkung}
		Damit wird die explizite Summation über alle Pfade operatorisch komprimiert. 
		Die Pfadsumme ist nun in den Potenzen eines Transferoperators enthalten. 
		Dies eröffnet den Zugang zu Spektralmethoden, dominanten Sektoren und asymptotischen Reduktionen.
	\end{bemerkung}
	
	\subsection{Diskrete Variationsrechnung}
	
	Die bisherige Konstruktion liefert eine Zustandssumme. Um ein Dynamikprinzip zu erhalten, wird nun eine diskrete Variationsrechnung eingeführt.
	
	\begin{definition}
		Sei
		\[
		\gamma=(X_0,\dots,X_N)
		\]
		ein diskreter BM-Pfad. 
		Eine \emph{lokale Variation} von \(\gamma\) an der Stelle \(k\in\{1,\dots,N-1\}\) ist die Ersetzung des inneren Zustands \(X_k\) durch einen Zustand \(Y\in\mathcal S\) mit
		\[
		X_{k-1}\rightsquigarrow Y\rightsquigarrow X_{k+1}.
		\]
		Der variierte Pfad sei
		\[
		\gamma^{(k\to Y)}
		=
		(X_0,\dots,X_{k-1},Y,X_{k+1},\dots,X_N).
		\]
	\end{definition}
	
	Da nur die beiden an \(X_k\) angrenzenden Kanten verändert werden, lautet die Wirkungsänderung
	\[
	\Delta_k(Y)
	=
	\mathcal S_{\mathrm{BM}}[\gamma^{(k\to Y)}]
	-
	\mathcal S_{\mathrm{BM}}[\gamma]
	\]
	explizit
	\[
	\Delta_k(Y)
	=
	L_{\mathrm{BM}}(X_{k-1},Y)
	+
	L_{\mathrm{BM}}(Y,X_{k+1})
	-
	L_{\mathrm{BM}}(X_{k-1},X_k)
	-
	L_{\mathrm{BM}}(X_k,X_{k+1}).
	\]
	
	\begin{definition}
		Ein diskreter BM-Pfad
		\[
		\gamma=(X_0,\dots,X_N)
		\]
		heißt \emph{lokal BM-stationär}, wenn für jeden inneren Index \(k\) der Zustand \(X_k\) ein lokaler Minimierer der Funktion
		\[
		Y\mapsto
		L_{\mathrm{BM}}(X_{k-1},Y)+L_{\mathrm{BM}}(Y,X_{k+1})
		\]
		über alle zulässigen \(Y\) ist.
	\end{definition}
	
	Dies ist die diskrete BM-Version einer Euler-Lagrange-Bedingung.
	
	\begin{satz}[diskrete BM-Euler-Lagrange-Bedingung]
		Ein diskreter BM-Pfad \(\gamma=(X_0,\dots,X_N)\) ist genau dann lokal BM-stationär, wenn für jeden inneren Index \(k\) gilt:
		\[
		X_k \in \operatorname*{arg\,min}_{Y:\,X_{k-1}\rightsquigarrow Y\rightsquigarrow X_{k+1}}
		\Bigl(
		L_{\mathrm{BM}}(X_{k-1},Y)+L_{\mathrm{BM}}(Y,X_{k+1})
		\Bigr).
		\]
	\end{satz}
	
	\begin{proof}
		Dies ist eine unmittelbare Umformulierung der lokalen Variationsbedingung \(\Delta_k(Y)\ge 0\) im Sinne des gewählten Ordnungs- beziehungsweise Minimierungsprinzips.
	\end{proof}
	
	\subsection{Reelle Zielfunktion als Alternativform}
	
	Da \(L_{\mathrm{BM}}\) komplexwertig ist, kann es für numerische oder modelltheoretische Zwecke nützlich sein, eine reelle Zielfunktion einzuführen. Eine mögliche Wahl ist
	\[
	\mathcal J[\gamma]
	=
	\alpha\,\Delta[\gamma]
	-
	\beta\,\mathrm{Coh}[\gamma],
	\]
	wobei
	\[
	\mathrm{Coh}[\gamma]
	=
	\left|
	\sum_{k=0}^{N-1}e^{i\phi(X_k,X_{k+1})}
	\right|
	\]
	ein Kohärenzmaß bezeichnet. Dann lautet das Variationsprinzip:
	\[
	\gamma_\star
	=
	\operatorname*{arg\,min}_{\gamma\in\Gamma(X,Y)}
	\mathcal J[\gamma].
	\]
	
	Diese Formulierung macht explizit, dass gute Pfade zugleich
	\begin{itemize}
		\item geringen kumulierten Defekt und
		\item hohe Phasenkohärenz
	\end{itemize}
	aufweisen sollen.
	
	\subsection{Spektrale Reduktion}
	
	Der Transferoperator eröffnet die Möglichkeit, die Forderung ``alle Pfade berücksichtigen'' analytisch zu entschärfen. 
	Statt alle Pfade explizit zu summieren, kann man das Spektrum von \(T\) untersuchen. 
	Falls \(T\) einen dominanten kohärenten Sektor besitzt, so ist für große \(N\) heuristisch
	\[
	T^N \approx T_{\mathrm{coh}}^N,
	\]
	wobei \(T_{\mathrm{coh}}\) auf einem defektarmen Unterraum operiert. 
	In diesem Sinn reduziert sich die volle Pfadsumme effektiv auf spektral dominante und strukturell verträgliche Klassen.
	
	\subsection{Interpretation}
	
	Die hier skizzierte Konstruktion verbindet drei Ebenen:
	
	\begin{enumerate}[label=\arabic*)]
		\item \textbf{diskreter Pfadraum:} Pfade sind Folgen von BM-Zuständen;
		\item \textbf{Variationsprinzip:} lokal stationäre Pfade minimieren eine diskrete BM-Wirkung;
		\item \textbf{operatorische Kompression:} der Transferoperator kodiert die gesamte Pfadsumme in seinen Potenzen.
	\end{enumerate}
	
	Damit entsteht ein erster analytischer Rahmen, in dem das Bamberger Modell eine Pfadintegral-ähnliche Dynamik formulieren kann, ohne auf eine buchstäbliche Summation über alle kontinuierlichen Bahnen angewiesen zu sein.
	
	\subsection{Ausblick}
	
	Die Lagrange-Vierlinge liefern in diesem Rahmen ein natürliches lokales Defektmodul. 
	Ein nächster Schritt wäre, konkrete BM-Zustandsräume zu fixieren und die Übergangsrelation \(\rightsquigarrow\) explizit aus den internen Regeln des Modells abzuleiten. 
	Darauf aufbauend könnten
	\begin{itemize}
		\item konkrete Transferoperatoren,
		\item spektrale Kohärenzsektoren,
		\item und numerische Minimierungsverfahren für BM-stationäre Pfade
	\end{itemize}
	systematisch untersucht werden.
	
\end{document}